%% paper_v1.tex
%% Explicit Evaluation of Sprung's Half-Logarithm H-Matrix
%% at X=0 for p=2 Supersingular Reduction
%%
%% Author: Raphael Elkuch
%% Date: 25 May 2026
%% Style: amsart, English, arXiv-ready

\documentclass[11pt,a4paper]{amsart}

\usepackage{amsmath,amssymb,amsthm,mathtools}
\usepackage[T1]{fontenc}
\usepackage{hyperref}
\hypersetup{colorlinks=true,linkcolor=blue,citecolor=blue,urlcolor=blue}

% Theorem environments
\newtheorem{theorem}{Theorem}[section]
\newtheorem{corollary}[theorem]{Corollary}
\newtheorem{lemma}[theorem]{Lemma}
\newtheorem{proposition}[theorem]{Proposition}
\theoremstyle{definition}
\newtheorem{definition}[theorem]{Definition}
\newtheorem{remark}[theorem]{Remark}
\newtheorem{example}[theorem]{Example}

% Macros
\newcommand{\Q}{\mathbb{Q}}
\newcommand{\Z}{\mathbb{Z}}
\newcommand{\F}{\mathbb{F}}
\newcommand{\C}{\mathbb{C}}
\newcommand{\Qp}{\mathbb{Q}_p}
\newcommand{\Zp}{\mathbb{Z}_p}
\newcommand{\Qtwo}{\mathbb{Q}_2}
\newcommand{\Ztwo}{\mathbb{Z}_2}
\newcommand{\eps}{\varepsilon}
\newcommand{\Lcal}{\mathcal{L}}
\newcommand{\Hcal}{\mathcal{H}}
\newcommand{\sharpflat}{{\sharp,\flat}}
\DeclareMathOperator{\Tr}{Tr}

\title[Sprung's $H$-matrix at $X=0$ for $p=2$ supersingular]{Explicit evaluation of Sprung's half-logarithm $H$-matrix at $X=0$ for $p=2$ supersingular reduction}

\author{Raphael Elkuch}
\email{relkuch@gmail.com}

\date{May 25, 2026}

\subjclass[2020]{Primary 11G05, 11G40, 11R23; Secondary 11Y40}

\keywords{$p$-adic $L$-function, supersingular reduction, Iwasawa theory, half-logarithm, Sprung's main conjecture, $p=2$ elliptic curves}

\begin{document}

\begin{abstract}
We give an explicit evaluation of Sprung's half-logarithm matrix $\Hcal(X)$ at $X=0$ for elliptic curves with supersingular reduction at $p=2$. The matrix limit defining $\Hcal(0)$ degenerates to the closed form $A^{-2}B$, where $A$ is the Frobenius action matrix and $B$ encodes the eigenvalues of the characteristic polynomial of Frobenius. As a consequence, the constant terms of the Sprung half-logarithms in their canonical basis are rational over $\Q$, although they a priori take values in the local quadratic extension $\Q_2(\alpha)$. We further give an explicit bridge $\eps^{-1}$ connecting Sage's Dieudonné basis to Sprung's canonical basis, and verify it bit-exactly on 180 elliptic curves drawn from the Cremona database with conductors in the range $[37,\,63\,096]$. The rationality phenomenon is interpreted as a renormalisation of the otherwise divergent Pollack half-logarithms at $T=0$, in the spirit of \cite[Remark 6.16]{Sprung2012}.
\end{abstract}

\maketitle

\section{Introduction}\label{sec:intro}

Iwasawa theory for elliptic curves at supersingular primes was reformulated by Kobayashi \cite{Kobayashi2003}, building on the construction of Pollack's plus/minus $p$-adic $L$-functions \cite{Pollack2003}, for the case $a_p = 0$. The case $a_p \neq 0$ was treated by Sprung \cite{Sprung2012}, who introduced a pair of half-logarithms $\log_\alpha^\sharp$, $\log_\alpha^\flat$ and a pair of $p$-adic $L$-functions $L_p^\sharp$, $L_p^\flat$ satisfying the canonical decomposition
\begin{equation}\label{eq:sprung-main}
L_p(E, \alpha, X) = L_p^\sharp(E, X) \cdot \log_\alpha^\sharp(X) + L_p^\flat(E, X) \cdot \log_\alpha^\flat(X).
\end{equation}
The matrix $\Hcal(X) = (\log_\alpha^\sharp, \log_\alpha^\flat)^T$ is defined as the projective limit
\begin{equation}\label{eq:H-def}
\Hcal(X) := \lim_{n \to \infty} C_1(X) \cdot C_2(X) \cdots C_n(X) \cdot A^{-(n+2)} \cdot B,
\end{equation}
where $C_k(X) = \begin{pmatrix} a_p & 1 \\ -\Phi_{p^k}(1+X) & 0 \end{pmatrix}$, $A$ is the Frobenius action matrix at $X=0$, and $B$ encodes the eigenvalues $\alpha, \beta$ of $X^2 - a_p X + p$. The convergence of this limit and its analytic properties have been studied in \cite{Sprung2012, Sprung2016}, but explicit closed-form evaluations are scarce.

The case $p = 2$ is delicate: Sprung's main theorem in \cite{Sprung2016} is stated for \emph{odd} supersingular primes, leaving the $p = 2$ case partly open in the literature. At the same time, computational software such as Sage's \texttt{padic\_lseries.Dp\_valued\_series} \cite{Sage} works with a slightly different basis convention, related to Sprung's by an $\eps$-bridge $\eps = (I - \varphi)^2$ where $\varphi$ is the Frobenius matrix in the Dieudonné basis $(\omega, \varphi\omega)$.

This note presents two observations.

\emph{(I) Triviality of the limit at $X=0$.} The defining limit \eqref{eq:H-def} degenerates algebraically at $X = 0$: since $\Phi_{p^k}(1) = p$ for all $k \geq 1$, the cofactor matrices $C_k(0)$ coincide with $A$, and $\Hcal(0) = A^{-2} B$ holds independently of the truncation level $n$. This is Theorem~\ref{thm:trivial-limit} below.

\emph{(II) The Sage--Sprung bridge.} The conversion between Sage's Dieudonné basis and Sprung's canonical basis is realised by an explicit rational matrix $\eps^{-1}$. We give the explicit form in Theorem~\ref{thm:bridge} and verify it bit-exactly on 180 elliptic curves.

These results are elementary, but useful as a working tool for computational Iwasawa theory at $p = 2$. They allow researchers to convert Sage's output to Sprung's convention without ambiguity, and they exhibit the constant components of Sprung's half-logarithms at $X = 0$ as elements of $\Q^2$, in spite of their a priori realisation in $\Q_2(\alpha)$.

\subsection*{Notation and conventions}
Throughout, $p = 2$ and $E/\Q$ is an elliptic curve without complex multiplication with good supersingular reduction at $p$, equivalently $a_p(E) \in \{-2, 0\}$. We write $\alpha, \beta$ for the eigenvalues of Frobenius acting on the $p$-adic Tate module, with characteristic polynomial $X^2 - a_p X + p$. We denote by $K_\alpha := \Q_2(\alpha)$ the local quadratic extension. The Frobenius matrix in the Dieudonné basis $(\omega, \varphi\omega)$ is taken to be
\begin{equation}
\varphi = \begin{pmatrix} 0 & -1/p \\ 1 & a_p/p \end{pmatrix},
\end{equation}
following the convention of Sage's implementation. The $\eps$-bridge is
\begin{equation}
\eps := (I - \varphi)^2 \in M_2(\Q).
\end{equation}


\section{Setting and preliminaries}\label{sec:setting}

\subsection{The Sprung half-logarithm matrix}\label{ssec:H-def}

Sprung's matrix $\Hcal(X)$ is defined by \eqref{eq:H-def}. We recall the auxiliary matrices: at $p = 2$,
\begin{align}
A &= \begin{pmatrix} a_p & 1 \\ -p & 0 \end{pmatrix} = \begin{pmatrix} a_p & 1 \\ -2 & 0 \end{pmatrix}, \\
B &= \begin{pmatrix} -1 & -1 \\ \beta & \alpha \end{pmatrix}.
\end{align}
The cyclotomic polynomials $\Phi_{p^k}(t)$ satisfy $\Phi_2(1) = 2$, $\Phi_4(1) = 2$, and more generally $\Phi_{2^k}(1) = 2$ for all $k \geq 1$ (the only prime dividing the order of a primitive root). Thus at $X = 0$:
\begin{equation}\label{eq:Phi-at-1}
\Phi_{p^k}(1) = p, \qquad k \geq 1.
\end{equation}

\subsection{The $\eps$-bridge}\label{ssec:eps-bridge}

In computational settings, $p$-adic $L$-functions and their decompositions \eqref{eq:sprung-main} are usually computed in Sage's Dieudonné basis $(\omega, \varphi\omega)$. The bridge to Sprung's canonical basis is given by $\eps = (I - \varphi)^2$.

\begin{lemma}\label{lem:det-eps}
$\det(\eps) = \dfrac{|E(\F_p)|^2}{p^2} = \dfrac{(p + 1 - a_p)^2}{p^2}$.
\end{lemma}

\begin{proof}
A direct computation gives
\[
\det(I - \varphi) = (1)\cdot\Bigl(1 - \frac{a_p}{p}\Bigr) - \Bigl(\frac{1}{p}\Bigr)\cdot(-1) = \frac{p + 1 - a_p}{p}.
\]
Squaring and applying $|E(\F_p)| = p + 1 - a_p$ yields the stated formula.
\end{proof}

In our setting $p = 2$, equation~\eqref{eq:det-eps-explicit} becomes:
\begin{equation}\label{eq:det-eps-explicit}
\det(\eps) = \frac{(3 - a_p)^2}{4} = \begin{cases} 25/4 & \text{if } a_p = -2, \\ 9/4 & \text{if } a_p = 0. \end{cases}
\end{equation}

\section{Main results}\label{sec:main}

\begin{theorem}[Triviality of the limit at $X = 0$]\label{thm:trivial-limit}
Let $p$ be any prime and let $A$ and $B$ be as in Section~\ref{ssec:H-def}. The Sprung half-logarithm matrix evaluated at $X = 0$ is
\begin{equation}\label{eq:H-zero}
\Hcal(0) = A^{-2} \cdot B,
\end{equation}
algebraically closed without nontrivial limit; in fact, for every $n \geq 1$ the truncated product satisfies
\begin{equation}\label{eq:Hn-zero}
C_1(0) \cdots C_n(0) \cdot A^{-(n+2)} = A^{-2}.
\end{equation}
\end{theorem}

\begin{proof}
By \eqref{eq:Phi-at-1}, $C_k(0) = A$ for every $k \geq 1$. Hence $C_1(0) \cdots C_n(0) = A^n$, and
\[
C_1(0) \cdots C_n(0) \cdot A^{-(n+2)} = A^n \cdot A^{-(n+2)} = A^{-2},
\]
independently of $n$. Passing to the limit and multiplying by $B$ yields $\Hcal(0) = A^{-2} B$.
\end{proof}

The matrix $A^{-2} \cdot B$ has explicit entries in $K_\alpha = \Q_2(\alpha)$. For $p = 2$:

\begin{proposition}[Explicit $\Hcal(0)$]\label{prop:H-explicit}
At $p = 2$ we have:
\begin{enumerate}
\item[(i)] For $a_p = -2$ ($\alpha = -1 + i$, $\beta = -1 - i$):
\begin{equation}
\Hcal(0) = \begin{pmatrix} -i/2 & i/2 \\ \tfrac{1}{2} - \tfrac{i}{2} & \tfrac{1}{2} + \tfrac{i}{2} \end{pmatrix}, \qquad \det \Hcal(0) = -\tfrac{i}{2}.
\end{equation}
\item[(ii)] For $a_p = 0$ ($\alpha = i\sqrt{2}$, $\beta = -i\sqrt{2}$):
\begin{equation}
\Hcal(0) = \begin{pmatrix} 1/2 & 1/2 \\ \tfrac{i\sqrt{2}}{2} & -\tfrac{i\sqrt{2}}{2} \end{pmatrix}, \qquad \det \Hcal(0) = -\tfrac{i\sqrt{2}}{2}.
\end{equation}
\end{enumerate}
\end{proposition}

\begin{proof}
For $a_p = -2$ one has $\det A = 2$ and
$A^{-2} = \begin{pmatrix} -1/2 & 1/2 \\ -1 & 1/2 \end{pmatrix}$.
The product $A^{-2} B$ with $B = \begin{pmatrix} -1 & -1 \\ -1-i & -1+i \end{pmatrix}$ yields the stated matrix. The case $a_p = 0$ is analogous.
\end{proof}

\begin{theorem}[Sage--Sprung bridge]\label{thm:bridge}
Let $E/\Q$ have $p = 2$ supersingular reduction. The conversion between Sage's basis $(\omega, \varphi\omega)$ and Sprung's canonical basis is the linear map
\begin{equation}
u_{\mathrm{Sprung}} = \eps^{-1} \cdot u_{\mathrm{Sage}}, \qquad \eps = (I - \varphi)^2.
\end{equation}
For the constant components of the Sprung half-logarithms in Sage's basis (with the normalisation of \cite[\S 5]{Sprung2012}), this yields the rational pairs:
\begin{enumerate}
\item[(i)] $a_p = -2$:
\begin{equation}
u_{\mathrm{Sage}} = (7/4,\, 3/4), \qquad u_{\mathrm{Sprung}} = (4/5,\, 9/10) \in \Q^2.
\end{equation}
\item[(ii)] $a_p = 0$:
\begin{equation}
v_{\mathrm{Sage}} = (1,\, 2), \qquad v_{\mathrm{Sprung}} = (-2/3,\, 4/3) \in \Q^2.
\end{equation}
\end{enumerate}
\end{theorem}

\begin{proof}
At $a_p = -2$ one computes
\[
\eps = \begin{pmatrix} 1/2 & 3/2 \\ -3 & 7/2 \end{pmatrix},
\qquad
\eps^{-1} = \begin{pmatrix} 14/25 & -6/25 \\ 12/25 & 2/25 \end{pmatrix},
\qquad
\det \eps = \tfrac{25}{4}.
\]
Direct multiplication gives $\eps^{-1} \cdot (7/4,\,3/4)^T = (4/5,\,9/10)^T$.
At $a_p = 0$ one similarly has
\[
\eps = \begin{pmatrix} 1/2 & 1 \\ -2 & 1/2 \end{pmatrix},
\qquad
\eps^{-1} = \begin{pmatrix} 2/9 & -4/9 \\ 8/9 & 2/9 \end{pmatrix},
\qquad
\det \eps = \tfrac{9}{4},
\]
and $\eps^{-1} \cdot (1,\,2)^T = (-2/3,\,4/3)^T$.
\end{proof}

\begin{corollary}[Rationality]\label{cor:rationality}
The constant components of Sprung's half-logarithms at $X = 0$ in the canonical basis are elements of $\Q^2$, although Sprung's matrix $\Hcal(0)$ takes a priori values in $K_\alpha = \Q_2(\alpha)$. Specifically:
\begin{equation}
u_{\mathrm{Sprung}}|_{a_p=-2} = (4/5,\, 9/10), \qquad v_{\mathrm{Sprung}}|_{a_p=0} = (-2/3,\, 4/3).
\end{equation}
\end{corollary}

\section{Numerical verification}\label{sec:numerical}

We verified Theorem~\ref{thm:bridge} on 180 elliptic curves drawn from the Cremona database \cite{Cremona} using SageMath \cite{Sage}. The sample is partitioned as follows: 75 curves with $a_p(E) = -2$ in the conductor range $[37,\,997]$, and 105 curves with $a_p(E) = 0$ in the conductor range $[77,\,63\,096]$. The $a_p = 0$ block is subdivided into 75 low-conductor curves in $[77,\,225]$ and 30 curves in $[1\,001,\,63\,096]$ obtained via logarithmically spaced bucket sampling. All curves have $\mathrm{rank}(E/\Q) \geq 1$.

For each curve we computed the Sage Dieudonné Frobenius matrix $\varphi$, formed $\eps = (I - \varphi)^2$, and applied $\eps^{-1}$ to $(7/4,\,3/4)^T$ if $a_p = -2$ or $(1,\,2)^T$ if $a_p = 0$. In every case the result equalled the predicted Sprung-basis vector $(4/5,\,9/10)^T$ or $(-2/3,\,4/3)^T$ respectively, bit-exactly in rational arithmetic over $\Q$.

\subsection{Sample of verified curves}

Table~\ref{tab:sample} shows a representative sample of 12 curves selected to span the full conductor range.

\begin{table}[h]
\centering
\begin{tabular}{l|r|r|r|c}
Label & Conductor & $a_p(E)$ & rank & Bridge match \\
\hline
\texttt{37a1}     & 37      & $-2$ & 1 & $\checkmark$ \\
\texttt{389a1}    & 389     & $-2$ & 2 & $\checkmark$ \\
\texttt{571b1}    & 571     & $-2$ & 2 & $\checkmark$ \\
\texttt{5077a1}   & 5\,077   & $-2$ & 3 & $\checkmark$ \\
\texttt{17963c1}  & 17\,963  & $-2$ & 2 & $\checkmark$ \\
\texttt{53461b1}  & 53\,461  & $-2$ & 3 & $\checkmark$ \\
\texttt{77a1}     & 77      & $0$  & 1 & $\checkmark$ \\
\texttt{1001a1}   & 1\,001   & $0$  & 1 & $\checkmark$ \\
\texttt{15855a1}  & 15\,855  & $0$  & 2 & $\checkmark$ \\
\texttt{30487a1}  & 30\,487  & $0$  & 3 & $\checkmark$ \\
\texttt{39812b1}  & 39\,812  & $0$  & 2 & $\checkmark$ \\
\texttt{63096b2}  & 63\,096  & $0$  & 1 & $\checkmark$ \\
\hline
\end{tabular}
\caption{Sample of curves verifying Theorem~\ref{thm:bridge}. Each row reports $\eps^{-1} \cdot u_{\mathrm{Sage}} = u_{\mathrm{Sprung}}$ bit-exactly in $\Q^2$.}
\label{tab:sample}
\end{table}

\subsection{Computational details}

The verification was conducted in SageMath~10 using the full Cremona database \cite{Cremona} (largest conductor $499\,998$). Stored rank information was used in lieu of \texttt{E.rank()} for performance. The total compute time was approximately five minutes on a standard workstation. The verification code and the complete log files are publicly available at
\begin{center}
\url{https://github.com/relkuch-del/sprung-h-matrix-p2}.
\end{center}

\subsection{Summary}

\begin{table}[h]
\centering
\begin{tabular}{l|r|l|r}
Class & $n$ & Conductor range & Bit-exact \\
\hline
$a_p = -2$ supersingular & 75 & $[37,\,997]$ & 75/75 \\
$a_p = 0$ supersingular (low)  & 75 & $[77,\,225]$ & 75/75 \\
$a_p = 0$ supersingular (high) & 30 & $[1\,001,\,63\,096]$ & 30/30 \\
\hline
\textbf{Total} & \textbf{180} & $[37,\,63\,096]$ & \textbf{180/180} \\
\end{tabular}
\caption{Summary of numerical verification. Zero mismatches and zero errors over four orders of magnitude in the conductor range.}
\label{tab:bilanz}
\end{table}

\section{Discussion: consistency with Pollack's half-logarithms}\label{sec:discussion}

In \cite[Remark 6.16]{Sprung2012}, Sprung relates his half-logarithms to Pollack's plus/minus logarithms $\log_2^\pm$ in the case $a_p = 0$:
\begin{equation}\label{eq:remark-6-16}
\log_\alpha^\sharp = -\tfrac{1}{2} \log_2^+ \cdot \alpha, \qquad \log_\alpha^\flat = \log_2^-.
\end{equation}
At $X = 0$, however, Pollack's logarithms behave singularly. Recall \cite[Lemma 4.1]{Pollack2003}:
\begin{equation}
\log_2^+(T) = \frac{1}{2} \prod_{n \geq 1} \frac{\Phi_{2n}(1+T)}{2}, \qquad \log_2^-(T) = \frac{1}{2} \prod_{n \geq 1} \frac{\Phi_{2n-1}(1+T)}{2}.
\end{equation}
Evaluating at $T = 0$:
\begin{itemize}
\item $\log_2^-(0) = 0$ \emph{structurally}, since $\Phi_1(1) = 0$ enters the product;
\item $\log_2^+(0)$ does \emph{not} converge in $\Q_2$: each truncation $N$ has 2-adic valuation $\leq -N$, with $v_2(\log_2^+(0)|_N) \to -\infty$ as $N \to \infty$.
\end{itemize}

Thus the constant values of Pollack's half-logarithms at $T = 0$ are not well-defined in $\Q_2$. The role of Sprung's normalisation factor $A^{-(n+2)}$ in \eqref{eq:H-def} is precisely to absorb this divergent behaviour into a well-defined element of $\Q^2$. Theorem~\ref{thm:trivial-limit} makes this concrete at $X = 0$: the renormalisation collapses to multiplication by $A^{-2}$, and the resulting matrix $A^{-2} B$ has rational entries in Sprung's canonical basis.

This structural consistency \emph{explains} the rationality observed in Theorem~\ref{thm:bridge} and Corollary~\ref{cor:rationality}: the vectors $u_{\mathrm{Sprung}}, v_{\mathrm{Sprung}} \in \Q^2$ are precisely the renormalised limits of the otherwise divergent Pollack constants.

\section{Open questions}\label{sec:open}

We close with three natural questions raised by, but not answered in, the present note.

\paragraph{Convergence of $\Hcal(X)$ for $X \neq 0$.}
The analysis here treats only the special point $X = 0$, where the limit degenerates algebraically. For $X \neq 0$ the cofactor matrices $C_k(X)$ depend nontrivially on $k$, and the convergence of the projective limit \eqref{eq:H-def} requires more delicate input. One expects an explicit description of $\Hcal(X)$ as a power series in $X$ over a suitable Banach algebra, extending the techniques of \cite{Sprung2012, Sprung2016}; this is left to future work.

\paragraph{The Iwasawa main conjecture at $p = 2$.}
The main theorem of \cite{Sprung2016} concerns odd supersingular primes. The case $p = 2$ is still open in the published literature. The bridge $\eps^{-1}$ provided by Theorem~\ref{thm:bridge} clarifies the relation between explicit Sage computations and Sprung's theoretical framework, and may be one of the technical ingredients for an eventual treatment of the $p = 2$ case.

\paragraph{Higher-order components.}
The vectors $u_{\mathrm{Sage}}, v_{\mathrm{Sage}}$ in Theorem~\ref{thm:bridge} are the constant components of Sage's $p$-adic $L$-series at $p = 2$. The full series expansions of $\log_\alpha^\sharp(X)$ and $\log_\alpha^\flat(X)$ around $X = 0$ display further rationality phenomena which warrant a separate investigation.

\section*{Acknowledgments}

The author thanks the SageMath developer community and J.~E.~Cremona for the elliptic-curve infrastructure on which this work depends. Manuscript preparation and bulk verification benefited from interactive work with Anthropic's Claude; all mathematical statements were independently verified by the author.

\begin{thebibliography}{99}

\bibitem{Kobayashi2003}
S.~Kobayashi,
\emph{Iwasawa theory for elliptic curves at supersingular primes,}
Invent. Math. \textbf{152} (2003), no.~1, 1--36.
DOI: \href{https://doi.org/10.1007/s00222-002-0265-4}{10.1007/s00222-002-0265-4}.

\bibitem{Pollack2003}
R.~Pollack,
\emph{On the $p$-adic $L$-function of a modular form at a supersingular prime,}
Duke Math.\ J.\ \textbf{118} (2003), no.~3, 523--558.
DOI: \href{https://doi.org/10.1215/S0012-7094-03-11835-9}{10.1215/S0012-7094-03-11835-9}.

\bibitem{Sprung2012}
F.~E.~I.~Sprung,
\emph{Iwasawa theory for elliptic curves at supersingular primes: A pair of main conjectures,}
J.\ Number Theory \textbf{132} (2012), no.~7, 1483--1506.
DOI: \href{https://doi.org/10.1016/j.jnt.2011.11.003}{10.1016/j.jnt.2011.11.003}.
arXiv: \href{https://arxiv.org/abs/0903.3419}{0903.3419}.

\bibitem{Sprung2016}
F.~E.~I.~Sprung,
\emph{The Iwasawa Main Conjecture for elliptic curves at odd supersingular primes,}
preprint, 2016.
arXiv: \href{https://arxiv.org/abs/1610.10017}{1610.10017}.

\bibitem{CastellaWan2016}
F.~Castella and X.~Wan,
\emph{Iwasawa Main Conjecture for Heegner Points: Supersingular Case,}
preprint, 2015/2016.
arXiv: \href{https://arxiv.org/abs/1506.02538}{1506.02538}.

\bibitem{Wan2014}
X.~Wan,
\emph{Iwasawa Main Conjecture for Supersingular Elliptic Curves and BSD Conjecture,}
preprint, 2014.
arXiv: \href{https://arxiv.org/abs/1411.6352}{1411.6352}.

\bibitem{Sage}
The Sage Developers,
\emph{SageMath, the Sage Mathematics Software System (Version 10.x),}
\url{https://www.sagemath.org}, 2026.

\bibitem{Cremona}
J.~E.~Cremona,
\emph{The Elliptic Curve Database for conductors up to $500\,000$,}
\url{https://github.com/JohnCremona/ecdata}, accessed May 2026.

\end{thebibliography}

\end{document}
